%% kapa: exam=1 round=35 year=2024 period=2 form=A subject=appraisal_law subject_ko=감정평가관계법규 seq=35 number=35
%% source: Q-Net 감정평가사 기출문제 (cst00309) artlSeq=5213229 file=raw/1008/5213229/2교시 문제지 원본.pdf
%% conversion: kordoc 4.24.51 -> tools/gen1.py (auto); review=proofread-sonnet; flags=none
\begin{question}{35}
부동산등기법령상 구분건물에 대한 등기에 관한 설명으로 옳지 않은 것은?
\begin{choices}
  \item 등기관은 구분건물인 경우 건물 등기기록의 표제부에 도면의 번호를 기록하여야 한다.
  \item 구분건물이 속하는 1동 전부가 멸실된 경우에는 그 구분건물의 소유권의 등기명의인은 1동의 건물에 속하는 다른 구분건물의 소유권의 등기명의인을 대위하여 1동 전부에 대한 멸실등기를 신청할 수 있다.
  \item 1동의 건물에 속하는 구분건물 중 일부만에 관하여 소유권보존등기를 신청하는 경우에는 나머지 구분건물의 표시에 관한 등기를 동시에 신청하여야 한다.
  \item 대지권이 등기된 구분건물의 등기기록에는 건물만에 관한 저당권설정등기를 할 수 있다.
  \item 구분건물에 대하여는 전유부분마다 부동산고유번호를 부여한다.
\end{choices}
\end{question}
\answer{4}
